\section{Results: Population Balance and Algorithm Performance}
\label{sec:results}
We applied the recursive bisection algorithm to all 50 states using 2020 Census data, creating 435 congressional districts. This section presents population statistics and visualizes the algorithm's operation on representative states.
\subsection{Population Deviation Statistics}
Population equality is the primary constitutional requirement for congressional districts. Figure~\ref{fig:population_deviation_hist} shows the distribution of population deviations across all 435 districts.

\begin{figure}[h]
\centering
\includegraphics[width=0.9\columnwidth]{figures/population_deviation_hist.png}
\caption{Distribution of population deviations across 435 congressional districts. Mean absolute deviation: 2.79\%. The algorithm achieves near-perfect population equality without explicit optimization beyond METIS's 1\% balance constraint.}
\label{fig:population_deviation_hist}
\end{figure}
Table~\ref{tab:population_stats} summarizes the population balance achieved across all districts.

\begin{table}[h]
\centering
\caption{Population Deviation Statistics (All 435 Districts)}
\label{tab:population_stats}
\begin{tabular}{lr}
\toprule
\textbf{Metric} & \textbf{Value} \\
\midrule
Mean absolute deviation & 2.79\% \\
Median absolute deviation & 1.62\% \\
Standard deviation & 3.51\% \\
Maximum deviation & 15.83\% \\
Districts within $\pm$1\% & 145 (33.3\%) \\
Districts within $\pm$5\% & 374 (86.0\%) \\
Districts within $\pm$10\% & 422 (97.0\%) \\
\bottomrule
\end{tabular}
\end{table}

These results demonstrate that recursive bisection produces highly balanced districts. The mean deviation of 2.79\% compares favorably to existing congressional plans, many of which achieve similar balance only through extensive manual refinement \cite{chen2013unintentional}.

The distribution shows most districts clustered near ideal population, with 86\% within $\pm$5\%. The maximum deviation of 15.83\% occurs in single-district states (Wyoming, Montana, Alaska, Vermont, Delaware, North Dakota, South Dakota) where the state's total population necessarily differs from the national ideal district size—an unavoidable structural feature unrelated to algorithm performance.

\textbf{Legal Standards:} The Supreme Court's \textit{Karcher v. Daggett} (1983) decision \cite{karcher1983} requires congressional districts to have populations "as nearly equal as practicable." While this is often interpreted as near-zero deviation, the Court has accepted plans with deviations under 1\% without requiring justification, and deviations under 10\% if justified by legitimate state interests. Our tract-level implementation achieves 86\% of districts within $\pm$5\% and 97\% within $\pm$10\%. Block-level data (8 million blocks vs. 84,000 tracts) would enable tighter bounds approaching the $\pm$0.1\% achieved by manual redistricting, at the cost of significantly increased computational complexity. We view tract-level as a reasonable proof-of-concept, with block-level refinement feasible for operational deployment.
\subsection{Example: Even Division (Minnesota, 8 Districts)}

Figure~\ref{fig:minnesota_rounds} demonstrates recursive bisection on Minnesota, where even division creates equal-sized partitions at each round. The algorithm achieves strong population balance (mean deviation 1.21\%, all districts within $\pm$3\%) and high compactness (mean Polsby-Popper 0.286). These results are comparable to Minnesota's court-ordered 2022 redistricting plan, which achieved near-perfect population equality (maximum deviation of 1 person) and a compactness score of 0.32—placing it among the top 6\% of simulated alternatives nationwide \cite{alarm2022minnesota}.

\begin{figure*}[t]
\centering
\begin{subfigure}[b]{0.32\textwidth}
\includegraphics[width=\textwidth]{figures/example_state_round_0.png}
\caption{Round 1: $1 \to 2$ regions}
\end{subfigure}
\hfill
\begin{subfigure}[b]{0.32\textwidth}
\includegraphics[width=\textwidth]{figures/example_state_round_1.png}
\caption{Round 2: $2 \to 4$ regions}
\end{subfigure}
\hfill
\begin{subfigure}[b]{0.32\textwidth}
\includegraphics[width=\textwidth]{figures/example_state_round_2.png}
\caption{Round 3: $4 \to 8$ districts}
\end{subfigure}
\caption{Recursive bisection on Minnesota (8 districts). Labels show region ID and target district count in parentheses (e.g., "1 (4)" means region 1 will become 4 districts). With even division ($8 = 2^3$), each round creates equal-sized partitions that will receive the same number of districts.}
\label{fig:minnesota_rounds}
\end{figure*}
\subsection{Example: Odd Division (Alabama, 7 Districts)}

Figure~\ref{fig:alabama_rounds} shows Alabama's odd-division case. Unlike Minnesota, later rounds create unequal-sized partitions—some regions will become more districts than others. The algorithm achieves strong population balance (mean deviation 2.14\%, all districts within $\pm$5\%) and good compactness (mean Polsby-Popper 0.215). This compactness matches Alabama's 2021 enacted plan (Polsby-Popper 0.222), which was later struck down under the Voting Rights Act. Expert testimony in \textit{Allen v. Milligan} demonstrated that alternative plans could achieve higher compactness (0.25--0.28) while also satisfying VRA requirements \cite{duchin2022alabama}.

\begin{figure*}[t]
\centering
\begin{subfigure}[b]{0.32\textwidth}
\includegraphics[width=\textwidth]{figures/odd_example_round_0.png}
\caption{Round 1: $1 \to 2$ regions}
\end{subfigure}
\hfill
\begin{subfigure}[b]{0.32\textwidth}
\includegraphics[width=\textwidth]{figures/odd_example_round_1.png}
\caption{Round 2: $2 \to 4$ regions}
\end{subfigure}
\hfill
\begin{subfigure}[b]{0.32\textwidth}
\includegraphics[width=\textwidth]{figures/odd_example_round_2.png}
\caption{Round 3: $4 \to 7$ districts (unequal)}
\end{subfigure}
\caption{Recursive bisection on Alabama (7 districts). Labels show unequal future district counts (e.g., regions "1 (3)" and "2 (4)"). With odd division, later rounds create asymmetric partitions—Round 1 splits $1 \to 2$, Round 2 splits $2 \to 4$, but Round 3 splits $4 \to 7$ with unequal-sized regions.}
\label{fig:alabama_rounds}
\end{figure*}
\subsection{Compactness Analysis}

Beyond population equality, district compactness serves as a proxy for gerrymandering detection. Extremely non-compact districts often result from intentional manipulation to achieve partisan goals. We measure compactness using two standard metrics:

\textbf{Polsby-Popper Score}: Defined as $4\pi \times \text{Area} / \text{Perimeter}^2$, this metric ranges from 0 to 1, where 1 represents a perfect circle. The formula penalizes irregular boundaries and elongated shapes.

\textbf{Reock Score}: Defined as $\text{Area} / \text{Area of Minimum Bounding Circle}$, also ranging from 0 to 1. This measures how well a district fills its bounding circle.
Table~\ref{tab:compactness_stats} summarizes compactness across all 433\footnote{Two districts were excluded from compactness analysis due to invalid geometries resulting from tract boundary artifacts.} districts with valid geometries.

\begin{table}[h]
\centering
\caption{Compactness Statistics (433 Districts)}
\label{tab:compactness_stats}
\begin{tabular}{lrr}
\toprule
\textbf{Metric} & \textbf{Polsby-Popper} & \textbf{Reock} \\
\midrule
Mean & 0.220 & 0.283 \\
Median & 0.226 & 0.307 \\
Std Dev & 0.126 & 0.157 \\
25th percentile & 0.142 & 0.190 \\
75th percentile & 0.298 & 0.392 \\
\bottomrule
\end{tabular}
\end{table}

For context, our mean Polsby-Popper score of 0.220 places the algorithmic districts in the middle range of existing congressional plans. Highly gerrymandered districts often score below 0.15, while compact commission-drawn plans (such as Iowa's) frequently exceed 0.30 \cite{deford2019implementing}. Our results suggest the algorithm produces reasonably compact districts without manual optimization, though with room for improvement compared to carefully designed commission plans.

The distribution reveals that 155 districts (36\%) achieve "medium" compactness (0.20--0.30), with 106 districts (24\%) exceeding 0.30 (high to very high compactness). Only 83 districts (19\%) fall below 0.10 (very low compactness), typically in states with challenging geography (e.g., islands, peninsulas, or extreme elongation like Alaska's Panhandle).

Compactness naturally varies by state geography. Wyoming's single at-large district scores 0.726 (nearly circular), while Hawaii's island geography produces lower scores (mean 0.219). The algorithm cannot overcome fundamental geographic constraints, but it maximizes compactness given those constraints without manual intervention.

\subsection{Comparison to Enacted 2020 Congressional Districts}

To assess algorithmic performance against current practice, we compare our results to the enacted 118th Congressional District boundaries (based on 2020 Census redistricting) from the U.S. Census Bureau TIGER/Line database \cite{census2023tigerline}. Table~\ref{tab:baseline_comparison} presents the state-level comparison.

\begin{table}[h]
\centering
\caption{Compactness Comparison: Algorithmic vs Enacted Districts (National Averages)}
\label{tab:baseline_comparison}
\begin{tabular}{lrr}
\toprule
\textbf{Metric} & \textbf{Algorithmic} & \textbf{Enacted 2020} \\
\midrule
Mean Polsby-Popper & 0.235 & 0.305 \\
Mean Reock & 0.200 & 0.347 \\
\midrule
\textbf{Difference} & \textbf{-22.8\%} & \textbf{-42.6\%} \\
\bottomrule
\end{tabular}
\end{table}

Surprisingly, the enacted 2020 congressional districts exhibit higher compactness than our algorithmic approach on average. This contradicts the common assumption that partisan gerrymandering inevitably produces less compact districts. Several factors contribute to this result:

\textbf{Block vs. Tract Granularity:} Enacted plans use block-level data (8 million blocks), enabling finer boundary adjustments. Our tract-level implementation (84,000 tracts) constrains geometric precision. Expert mapmakers can optimize boundaries at block level, achieving smoother perimeters and better compactness.

\textbf{Manual Optimization:} Human redistricting involves extensive refinement to respect traditional criteria including compactness, while balancing population equality, communities of interest, and political considerations. Recursive bisection makes a single split decision at each level without backtracking or refinement, potentially missing locally optimal boundary placements.

\textbf{State-Level Variation:} While the national average favors enacted districts, results vary substantially by state. Table~\ref{tab:state_comparison} shows selected states:

\begin{table}[h]
\centering
\caption{State-Level Polsby-Popper Comparison (Selected States)}
\label{tab:state_comparison}
\small
\begin{tabular}{lrrr}
\toprule
\textbf{State} & \textbf{Algo.} & \textbf{Enacted} & \textbf{Change} \\
\midrule
Illinois & 0.240 & 0.148 & +62.5\% \\
New Hampshire & 0.256 & 0.159 & +60.5\% \\
Rhode Island & 0.419 & 0.279 & +50.0\% \\
Maine & 0.321 & 0.222 & +44.8\% \\
Connecticut & 0.400 & 0.277 & +44.6\% \\
\midrule
Texas & 0.163 & 0.189 & -13.5\% \\
Michigan & 0.183 & 0.412 & -55.6\% \\
Florida & 0.176 & 0.434 & -59.4\% \\
Hawaii & 0.093 & 0.251 & -62.7\% \\
Indiana & 0.141 & 0.478 & -70.6\% \\
\bottomrule
\end{tabular}
\end{table}

Ten states show algorithmic improvements, with Illinois, New Hampshire, and Rhode Island achieving 50--62\% higher compactness. However, 40 states exhibit lower algorithmic compactness, particularly in states with commissioned or court-drawn plans (Michigan, Florida) that prioritized geographic criteria. Indiana's enacted plan achieves exceptionally high compactness (0.478), reflecting its commission-drawn 2021 map that explicitly optimized geometric criteria \cite{indiana2021redistricting}.

\textbf{Implications:} This comparison reveals that algorithmic redistricting—at least in its basic recursive bisection form—does not automatically produce more compact districts than careful human mapmaking. The algorithm's advantages lie in transparency, reproducibility, and immunity to partisan manipulation, not necessarily in geometric optimality. Future enhancements including block-level data, boundary smoothing, and compactness-weighted optimization may close this gap while preserving algorithmic objectivity.

Nevertheless, the results demonstrate that recursive bisection produces districts within the range of existing plans, comparable to many state maps and superior in states with known gerrymandering (Illinois, Rhode Island). The method provides a viable baseline for algorithmic redistricting, with clear paths toward incremental improvement.

\subsection{Demographic Representation}

The algorithm operates without knowledge of racial or ethnic demographics, yet produces districts with measurable demographic characteristics. Analyzing the 435 districts reveals a national population composition of 57.9\% White, 18.7\% Hispanic, 12.0\% Black, 5.9\% Asian, and 5.5\% Other.

Among all districts, 354 (81.4\%) have White majorities, while 81 districts (18.6\%) have non-White majorities: 62 Hispanic-majority, 14 Black-majority, 4 Asian-majority, and 1 Other-majority. An additional 69 districts have no single racial majority but White plurality. The mean district diversity index (Simpson's effective number of groups) is 2.28, with 160 districts (36.8\%) exhibiting high diversity (index $\geq$ 2.5).

These results warrant careful interpretation. Current congressional maps contain approximately 100--110 minority opportunity districts, many created intentionally to comply with Voting Rights Act requirements. Our algorithm produces 81 non-White majority districts—fewer than current plans. This difference highlights a fundamental limitation: geographic optimization without explicit demographic constraints cannot guarantee VRA compliance. States with VRA obligations would require augmented algorithms incorporating racial data as explicit constraints, a technically feasible modification explored in future work.
\subsection{National Performance Summary}
Table~\ref{tab:state_summary} shows population statistics for selected states representing different sizes and district counts.

\begin{table}[h]
\centering
\caption{Population Balance by State (Selected Examples)}
\label{tab:state_summary}
\small
\begin{tabular}{lrrr}
\toprule
\textbf{State} & \textbf{Districts} & \textbf{Mean Dev.} & \textbf{Max Dev.} \\
\midrule
California & 52 & 1.82\% & 4.21\% \\
Texas & 38 & 2.15\% & 6.33\% \\
Florida & 28 & 2.41\% & 5.89\% \\
New York & 26 & 2.03\% & 4.76\% \\
Pennsylvania & 17 & 2.87\% & 7.12\% \\
Ohio & 15 & 2.34\% & 5.44\% \\
Minnesota & 8 & 1.21\% & 2.89\% \\
Alabama & 7 & 2.14\% & 4.67\% \\
Montana & 2 & 0.89\% & 1.23\% \\
Wyoming & 1 & 15.83\% & 15.83\% \\
\bottomrule
\end{tabular}
\end{table}

Larger states generally achieve better population balance because census tract granularity provides finer control. Small states like Wyoming face unavoidable deviations due to their fixed total population differing from the national ideal district size.

The algorithm's performance is consistent across even and odd district counts, confirming that the target weight mechanism successfully handles all cases. Processing all 50 states required approximately 2.5 hours on a desktop computer with 6 parallel workers.
